\let\negmedspace\undefined
\let\negthickspace\undefined
\documentclass[journal]{IEEEtran}
\usepackage[a5paper, margin=10mm, onecolumn]{geometry}
\usepackage{tfrupee} % Include tfrupee package

\setlength{\headheight}{1cm} 
\setlength{\headsep}{0mm}     

\usepackage{gvv-book}
\usepackage{gvv}
\usepackage{cite}
\usepackage{amsmath,amssymb,amsfonts,amsthm}
\usepackage{algorithmic}
\usepackage{graphicx}
\usepackage{textcomp}
\usepackage{xcolor}
\usepackage{txfonts}
\usepackage{listings}
\usepackage{enumitem}
\usepackage{mathtools}
\usepackage{gensymb}
\usepackage{comment}
\usepackage[breaklinks=true]{hyperref}
\usepackage{tkz-euclide} 
\usepackage{listings}
\def\inputGnumericTable{}                                 
\usepackage[latin1]{inputenc}                                
\usepackage{color}                                            
\usepackage{array}                                            
\usepackage{longtable}                                       
\usepackage{calc}                                             
\usepackage{multirow}                                         
\usepackage{hhline}                                           
\usepackage{ifthen}                                           
\usepackage{lscape}
\usepackage{circuitikz}
\tikzstyle{block} = [rectangle, draw, fill=blue!20, 
    text width=4em, text centered, rounded corners, minimum height=3em]
\tikzstyle{sum} = [draw, fill=blue!10, circle, minimum size=1cm, node distance=1.5cm]
\tikzstyle{input} = [coordinate]
\tikzstyle{output} = [coordinate]


\begin{document}

\bibliographystyle{IEEEtran}
\vspace{3cm}

\title{5.13.12}
\author{AI25BTECH11036-SNEHAMRUDULA}
 \maketitle
{\let\newpage\relax\maketitle}
\renewcommand{\thefigure}{\theenumi}
\renewcommand{\thetable}{\theenumi}
\setlength{\intextsep}{10pt} 
\numberwithin{equation}{enumi}
\numberwithin{figure}{enumi}
\renewcommand{\thetable}{\theenumi}
\textbf{Question}:\\
The set of all values of $\lambda$ for which the system of linear equations
\[
\begin{aligned}
2x_{1} - 2x_{2} + x_{3} &= \lambda x_{1}, \\[6pt]
2x_{1} - 3x_{2} + 2x_{3} &= \lambda x_{2}, \\[6pt]
- x_{1} + 2x_{2} &= \lambda x_{3},
\end{aligned}
\]
has a non-trivial solution
\begin{enumerate}
  \item[(a)] contains two elements
  \item[(b)] contains more than two elements
  \item[(c)] is an empty set
  \item[(d)] is a singleton
\end{enumerate}

\solution \\
%-----------------------------------------
% Q 3.4.5 : Rhombus construction (gvv methods)
%-----------------------------------------
Consider the coefficient matrix
\begin{align}
A &= \myvec{2 & -2 & 1 \\ 2 & -3 & 2 \\ -1 & 2 & 0}.
\end{align}

The characteristic equation is
\begin{align}
\det(A - \lambda I) &= 0.
\end{align}

On solving, the eigenvalues are obtained as
\begin{align}
\lambda_1 &= -3, \quad \lambda_2 = 1, \quad \lambda_3 = 1.
\end{align}

Corresponding eigenvectors are
\begin{align}
v_1 &= \myvec{-1 \\ -2 \\ 1}, \quad
v_2 = \myvec{2 \\ 1 \\ 0}, \quad
v_3 = \myvec{-1 \\ 0 \\ 1}.
\end{align}

Thus, the modal matrix is
\begin{align}
P &= \myvec{-1 & 2 & -1 \\ -2 & 1 & 0 \\ 1 & 0 & 1}.
\end{align}

The diagonal eigenvalue matrix is
\begin{align}
D &= \myvec{-3 & 0 & 0 \\ 0 & 1 & 0 \\ 0 & 0 & 1}.
\end{align}

Inverse of $P$ is
\begin{align}
P^{-1} &= \myvec{\tfrac{1}{4} & -\tfrac{1}{2} & \tfrac{1}{4} \\[6pt]
\tfrac{1}{2} & 0 & \tfrac{1}{2} \\[6pt]
-\tfrac{1}{4} & \tfrac{1}{2} & \tfrac{3}{4}}.
\end{align}

Hence, the eigendecomposition is
\begin{align}
A &= P D P^{-1}.
\end{align}

\begin{figure}[h!]
    \centering
    \includegraphics[height=0.4\textheight, keepaspectratio]{5.13.12fig.png}

      \includegraphics[height=0.4\textheight, keepaspectratio]{5.13.12fig(1).png}
   
    \label{figure_1}
\end{figure}
 
\end{document}
